\section{Conclusion}
Within these passive static comparisons, sex alone was an incomplete descriptor of sacroiliac response. Explicit pelvic size and dimensions substantially attenuated sex-associated translation differences, while individual geometry retained much of the modeled response under a standardized bone-density field and fixed soft-tissue properties. The resulting motion depended on load distribution and application site. Pelvic function is therefore more informatively described through individual architecture interacting with loading than through a fixed sex-specific mobility value; this interpretation does not establish sex equivalence, causal mediation or physiological validation.
